\documentclass[a4paper,11pt]{article}

\usepackage[a4paper,margin=2cm]{geometry}
\usepackage[T1]{fontenc}
\usepackage[utf8]{inputenc}
\usepackage[ngerman,english]{babel}
\usepackage{lmodern}
\usepackage{microtype}
\usepackage{xcolor}
\usepackage{parskip}
\usepackage{enumitem}
\usepackage[most]{tcolorbox}
\usepackage{titlesec}
\usepackage[hidelinks]{hyperref}
\usepackage{array}
\usepackage{longtable}
\usepackage[table]{xcolor}

\definecolor{HeaderBlueDark}{RGB}{70,110,170}
\definecolor{RowBlueLight}{RGB}{235,243,255}
\definecolor{RowBlueDark}{RGB}{220,233,250}
\definecolor{SoftGray}{RGB}{90,90,90}

\setlength{\parindent}{0pt}
\setlist[itemize]{leftmargin=1.4em,itemsep=0.35em,topsep=0.35em}
\linespread{1.00}
\renewcommand{\arraystretch}{1.35}
\setlength{\tabcolsep}{6pt}

\titleformat{\section}[block]
{\Large\bfseries\color{HeaderBlueDark}}
{}{0pt}{}
\titlespacing{\section}{0pt}{1.0em}{0.6em}

\titleformat{\subsection}[block]
{\large\bfseries\color{HeaderBlueDark}}
{}{0pt}{}
\titlespacing{\subsection}{0pt}{0.8em}{0.5em}

\newtcolorbox{exbox}{enhanced,breakable,colback=RowBlueLight,colframe=RowBlueDark,boxrule=0.4pt,arc=1.5mm,left=1.2mm,right=1.2mm,top=1mm,bottom=1mm,title={Beispiele \textnormal{(Examples)}},fonttitle=\bfseries,colbacktitle=RowBlueDark,coltitle=black}

\NewDocumentEnvironment{grammarentry}{mm+b}{%
    \begin{tcolorbox}[enhanced,breakable,colback=white,colframe=HeaderBlueDark,boxrule=0.6pt,arc=1.5mm,left=1.5mm,right=1.5mm,top=1mm,bottom=1mm,colbacktitle=HeaderBlueDark,coltitle=white,fonttitle=\bfseries,title={#1\hfill\normalfont\footnotesize #2}]
    #3
    \end{tcolorbox}
}{}

\newcommand{\GermanText}[1]{\textbf{Deutsch: }#1\par}
\newcommand{\EnglishText}[1]{{\small\color{SoftGray}\textbf{English: }#1\par}}
\newcommand{\ExampleItem}[2]{\item \textbf{#1}\\{\small\color{SoftGray}#2}}

\title{Das deutsche Verb \glqq stammen\grqq\\\large The German Verb ``stammen''}
\date{}

\begin{document}
\maketitle
\tableofcontents
\newpage

\section{Grundbedeutung \textnormal{(Basic meaning)}}
\begin{grammarentry}{stammen}{Verb, intransitiv, regelmäßig}
\GermanText{\textbf{stammen} bedeutet, seinen Ursprung oder seine Herkunft in einer Person, einem Ort, einer Zeit oder einer Quelle zu haben.}
\EnglishText{\textbf{stammen} means to have one's origin or provenance in a person, place, time, or source.}
\GermanText{Aussprache: Betonung auf der ersten Silbe: \textbf{STAM-men}.}
\EnglishText{Pronunciation: approximately \textit{SHTAHM-men}.}
\GermanText{Das Verb ist regelmäßig und nicht trennbar. Es wird normalerweise ohne Akkusativobjekt gebraucht.}
\EnglishText{The verb is regular and non-separable. It is normally used without an accusative object.}
\end{grammarentry}

\section{Die wichtigste Struktur: \glqq aus + Dativ stammen\grqq\ \textnormal{(The most important structure: ``stammen aus + dative'')}}
\begin{grammarentry}{aus + Dativ}{Herkunft und Ursprung}
\GermanText{Mit \textbf{aus + Dativ} bezeichnet \textbf{stammen} meistens die geografische, zeitliche, materielle oder soziale Herkunft.}
\EnglishText{With \textbf{aus + dative}, \textbf{stammen} usually expresses geographical, temporal, material, or social origin.}
\GermanText{Struktur: \textbf{jemand/etwas stammt aus + Dativ}.}
\EnglishText{Structure: \textbf{someone/something comes from + dative}.}
\end{grammarentry}
\begin{exbox}
\begin{itemize}
\ExampleItem{Ich stamme aus dem Iran.}{I come from Iran.}
\ExampleItem{Sie stammt aus einer kleinen Stadt.}{She comes from a small town.}
\ExampleItem{Dieses Möbelstück stammt aus dem 19. Jahrhundert.}{This piece of furniture dates from the 19th century.}
\ExampleItem{Das Wort stammt aus dem Lateinischen.}{The word comes from Latin.}
\end{itemize}
\end{exbox}

\section{Die Struktur \glqq von + Dativ stammen\grqq\ \textnormal{(The structure ``stammen von + dative'')}}
\begin{grammarentry}{von + Dativ}{Urheber, Quelle oder Herkunft}
\GermanText{Mit \textbf{von + Dativ} nennt man häufig eine Person oder Quelle, von der etwas kommt oder geschaffen wurde.}
\EnglishText{With \textbf{von + dative}, one often names a person or source from which something comes or by whom it was created.}
\GermanText{Diese Struktur ist besonders häufig bei Zitaten, Ideen, Werken, Informationen und Gegenständen.}
\EnglishText{This structure is especially common with quotations, ideas, works, information, and objects.}
\end{grammarentry}
\begin{exbox}
\begin{itemize}
\ExampleItem{Dieses Zitat stammt von Goethe.}{This quotation is by Goethe.}
\ExampleItem{Die Idee stammt von meiner Kollegin.}{The idea comes from my colleague.}
\ExampleItem{Die Information stammt von einer zuverlässigen Quelle.}{The information comes from a reliable source.}
\ExampleItem{Der Brief stammt von meinem Großvater.}{The letter is from my grandfather.}
\end{itemize}
\end{exbox}

\section{\glqq aus\grqq\ oder \glqq von\grqq? \textnormal{(``aus'' or ``von''?)}}
\rowcolors{2}{RowBlueLight}{RowBlueDark}
\begin{longtable}{>{\bfseries}p{3.0cm}|p{5.2cm}|p{6.0cm}}
\rowcolor{HeaderBlueDark}
\color{white}\textbf{Struktur / Structure} & \color{white}\textbf{Deutsch / German} & \color{white}\textbf{English} \\
\textbf{aus + Dativ} & Herkunft aus einem Ort, einer Zeit, Sprache, Gruppe oder einem Bereich. & Origin from a place, period, language, group, or domain. \\
\textbf{von + Dativ} & Ursprung bei einer Person oder konkreten Quelle; oft Urheber oder Absender. & Origin from a person or concrete source; often author, creator, or sender. \\
\end{longtable}
\begin{exbox}
\begin{itemize}
\ExampleItem{Das Gemälde stammt aus dem 18. Jahrhundert.}{The painting dates from the 18th century.}
\ExampleItem{Das Gemälde stammt von einem österreichischen Künstler.}{The painting is by an Austrian artist.}
\end{itemize}
\end{exbox}

\section{Personen und Herkunft \textnormal{(People and origin)}}
\begin{grammarentry}{stammen aus}{bei Personen}
\GermanText{Bei Personen bedeutet \textbf{aus einem Land, einer Stadt oder einer Familie stammen}, dort seine Herkunft zu haben.}
\EnglishText{With people, \textbf{to come from a country, city, or family} means to have one's origin there.}
\GermanText{\textbf{Ich stamme aus Österreich} beschreibt die Herkunft. Für den heutigen Wohnort benutzt man dagegen normalerweise \textbf{wohnen} oder \textbf{leben}.}
\EnglishText{\textbf{Ich stamme aus Österreich} describes origin. For one's current place of residence, German normally uses \textbf{wohnen} or \textbf{leben}.}
\end{grammarentry}
\begin{exbox}
\begin{itemize}
\ExampleItem{Er stammt aus Wien, lebt aber jetzt in Graz.}{He comes from Vienna but now lives in Graz.}
\ExampleItem{Sie stammt aus einer Musikerfamilie.}{She comes from a family of musicians.}
\end{itemize}
\end{exbox}

\section{Unterschied zwischen \glqq stammen\grqq\ und \glqq kommen\grqq\ \textnormal{(Difference between ``stammen'' and ``kommen'')}}
\begin{grammarentry}{stammen vs. kommen}{Herkunft}
\GermanText{\textbf{kommen aus} und \textbf{stammen aus} können beide Herkunft ausdrücken. Bei Personen ist \textbf{kommen aus} im Alltag meist neutraler und häufiger. \textbf{stammen aus} klingt oft etwas sachlicher und betont stärker die Herkunft oder den Ursprung.}
\EnglishText{\textbf{kommen aus} and \textbf{stammen aus} can both express origin. With people, \textbf{kommen aus} is usually more neutral and more common in everyday speech. \textbf{stammen aus} often sounds somewhat more factual and emphasizes origin more strongly.}
\GermanText{Bei historischen Gegenständen, Wörtern, Ideen und Quellen ist \textbf{stammen} besonders typisch.}
\EnglishText{With historical objects, words, ideas, and sources, \textbf{stammen} is especially typical.}
\end{grammarentry}
\begin{exbox}
\begin{itemize}
\ExampleItem{Ich komme aus Graz.}{I come from Graz.}
\ExampleItem{Ich stamme aus Graz.}{I originate from / come from Graz.}
\ExampleItem{Dieses Dokument stammt aus dem Jahr 1920.}{This document dates from the year 1920.}
\end{itemize}
\end{exbox}

\section{Unterschied zwischen \glqq stammen\grqq\ und \glqq abstammen\grqq\ \textnormal{(Difference between ``stammen'' and ``abstammen'')}}
\begin{grammarentry}{abstammen von}{Abstammung und Vorfahren}
\GermanText{\textbf{abstammen von + Dativ} benutzt man besonders für biologische oder genealogische Abstammung, also für Vorfahren und Entwicklungslinien.}
\EnglishText{\textbf{abstammen von + dative} is used especially for biological or genealogical descent, that is, for ancestors and lines of development.}
\GermanText{\textbf{stammen aus} ist allgemeiner und beschreibt Herkunft; \textbf{abstammen von} betont Abstammung von Vorfahren oder einer früheren Form.}
\EnglishText{\textbf{stammen aus} is more general and describes origin; \textbf{abstammen von} emphasizes descent from ancestors or an earlier form.}
\end{grammarentry}
\begin{exbox}
\begin{itemize}
\ExampleItem{Er stammt aus einer alten Familie.}{He comes from an old family.}
\ExampleItem{Er stammt von einer alten Familie ab.}{He descends from an old family.}
\end{itemize}
\end{exbox}

\section{Konjugation \textnormal{(Conjugation)}}
\rowcolors{2}{RowBlueLight}{RowBlueDark}
\begin{longtable}{>{\bfseries}p{3.4cm}|p{5.0cm}|p{5.5cm}}
\rowcolor{HeaderBlueDark}
\color{white}\textbf{Form / Form} & \color{white}\textbf{Deutsch / German} & \color{white}\textbf{English} \\
Präsens & ich stamme & I come from / originate \\
Präteritum & ich stammte & I came from / originated \\
Perfekt & ich habe gestammt & I have come from / originated \\
Plusquamperfekt & ich hatte gestammt & I had come from / originated \\
Futur I & ich werde stammen & I will come from / originate \\
Konjunktiv II & ich würde stammen & I would come from / originate \\
Konjunktiv I & ich stamme & I come from / originate \\
Partizip II & gestammt & originated / come from \\
\end{longtable}
\begin{grammarentry}{Hinweis zum Perfekt}{usage note}
\GermanText{Die Perfektform \textbf{hat gestammt} ist grammatisch möglich, aber in vielen Herkunftsaussagen verwendet man häufiger Präsens oder Präteritum, weil die Herkunft meist als dauerhafte Tatsache verstanden wird.}
\EnglishText{The perfect form \textbf{hat gestammt} is grammatically possible, but in many statements of origin, the present or simple past is more common because origin is usually understood as a lasting fact.}
\end{grammarentry}

\section{Typische Bedeutungsbereiche \textnormal{(Typical semantic areas)}}
\begin{grammarentry}{Herkunft}{wichtige Verwendungen}
\GermanText{\textbf{Person:} aus einem Land, einer Stadt, Region oder Familie stammen.}
\EnglishText{\textbf{Person:} to come from a country, city, region, or family.}
\GermanText{\textbf{Zeit:} aus einer bestimmten Epoche oder einem Jahr stammen.}
\EnglishText{\textbf{Time:} to date from a particular period or year.}
\GermanText{\textbf{Sprache:} ein Wort stammt aus einer anderen Sprache.}
\EnglishText{\textbf{Language:} a word comes from another language.}
\GermanText{\textbf{Urheber:} ein Zitat, eine Idee oder ein Werk stammt von jemandem.}
\EnglishText{\textbf{Creator/source:} a quotation, idea, or work comes from / is by someone.}
\GermanText{\textbf{Quelle:} Daten oder Informationen stammen aus einer Studie oder von einer Institution.}
\EnglishText{\textbf{Source:} data or information comes from a study or institution.}
\end{grammarentry}

\section{Typische Fehler \textnormal{(Common mistakes)}}
\begin{grammarentry}{Kasus und Präposition}{wichtig}
\GermanText{Nach \textbf{aus} und \textbf{von} steht immer der \textbf{Dativ}.}
\EnglishText{After \textbf{aus} and \textbf{von}, German always uses the \textbf{dative}.}
\GermanText{Richtig: \textbf{Ich stamme aus dem Iran.} Nicht: \textbf{aus den Iran}.}
\EnglishText{Correct: \textbf{Ich stamme aus dem Iran.} Not: \textbf{aus den Iran}.}
\GermanText{Richtig: \textbf{Die Idee stammt von meiner Freundin.} Nicht: \textbf{von meine Freundin}.}
\EnglishText{Correct: \textbf{Die Idee stammt von meiner Freundin.} Not: \textbf{von meine Freundin}.}
\end{grammarentry}

\section{Zusammenfassung \textnormal{(Summary)}}
\begin{grammarentry}{Merksatz}{kurz und wichtig}
\GermanText{\textbf{stammen aus + Dativ} = Herkunft aus einem Ort, einer Zeit, Sprache, Familie oder einem Bereich.}
\EnglishText{\textbf{stammen aus + dative} = origin from a place, time, language, family, or domain.}
\GermanText{\textbf{stammen von + Dativ} = Ursprung bei einer Person oder konkreten Quelle; oft Urheber, Absender oder Quelle.}
\EnglishText{\textbf{stammen von + dative} = origin from a person or concrete source; often author, sender, or source.}
\GermanText{\textbf{abstammen von + Dativ} = genealogisch oder biologisch von Vorfahren bzw. einer früheren Form herkommen.}
\EnglishText{\textbf{abstammen von + dative} = to descend genealogically or biologically from ancestors or an earlier form.}
\end{grammarentry}

\end{document}
